% ═════════════════════════════════════════════════════════════════════════════
%  chapitres/introduction.tex
% ═════════════════════════════════════════════════════════════════════════════
\chapter*{Introduction Générale}
\addcontentsline{toc}{chapter}{Introduction Générale}
\markboth{Introduction Générale}{Introduction Générale}

\section*{Contexte et problématique}

L'industrie du logiciel connaît depuis une décennie une transformation profonde,
portée par deux évolutions majeures. D'une part, le \textbf{passage aux architectures
microservices} a fragmenté les applications monolithiques en un ensemble de services
indépendants, chacun pouvant être développé, déployé et mis à l'échelle
individuellement. D'autre part, l'adoption généralisée des pratiques
\textbf{DevOps} a raccourci drastiquement les cycles de livraison, permettant de
passer de plusieurs déploiements par an à plusieurs déploiements par jour.

Si ces deux tendances ont considérablement amélioré l'agilité et la résilience des
systèmes, elles ont simultanément introduit une \textbf{surface d'attaque étendue}
et difficile à maîtriser. Un système composé de dizaines, voire de centaines de
microservices communicants génère autant de canaux de communication potentiellement
vulnérables, et la vitesse de déploiement laisse peu de place aux audits de sécurité
traditionnels, souvent relégués en fin de cycle.

\begin{importantbox}
  \textbf{Constat fondamental :} Selon le rapport \textit{State of DevOps} 2023,
  plus de 60\,\% des incidents de sécurité en production dans les environnements
  cloud-natifs sont liés à des communications interservices non sécurisées ou à
  des configurations d'accès insuffisamment contrôlées.
\end{importantbox}

Face à ce constat, la communauté technique a progressivement convergé vers deux
approches complémentaires. Le paradigme \textbf{DevSecOps} préconise d'intégrer
la sécurité dès les premières phases du développement --- le principe du
\textit{Shift Left Security} --- en automatisant les contrôles de sécurité au sein
des pipelines CI/CD. Le \textbf{Service Mesh}, quant à lui, propose une couche
d'infrastructure dédiée à la gestion des communications interservices, offrant
nativement des mécanismes de chiffrement, d'authentification et d'observabilité
sans modifier le code applicatif.

La problématique centrale de ce projet est donc la suivante :

\begin{tcolorbox}[
  colback=gray!6, colframe=darkgray,
  fonttitle=\bfseries, title={Question de recherche},
  arc=4pt
]
  \textit{Dans quelle mesure l'intégration d'un Service Mesh au sein d'une
  démarche DevSecOps permet-elle de renforcer significativement la sécurité
  des communications interservices dans les architectures microservices, et
  quelles sont les conditions d'une telle intégration ?}
\end{tcolorbox}

\section*{Objectifs du projet}

Ce projet de fin d'année poursuit quatre objectifs principaux :

\begin{enumerate}
  \item \textbf{Analyser} les pratiques DevSecOps et les différents types de
    tests de sécurité applicables aux architectures microservices, en mettant
    en évidence les limites des approches conventionnelles.

  \item \textbf{Étudier} le concept de Service Mesh, son architecture, ses
    composants fondamentaux et le panorama des solutions existantes (Istio,
    Linkerd, Consul, Cilium), afin d'identifier la solution la mieux adaptée
    à nos contraintes.

  \item \textbf{Modéliser} les menaces pesant sur une architecture microservices,
    selon la méthodologie STRIDE, et comparer les profils de risque avec et sans
    Service Mesh.

  \item \textbf{Concevoir et valider} un prototype fonctionnel sur Kubernetes avec
    Istio, intégrant les mécanismes de sécurité mTLS, JWT et RBAC, et en mesurer
    l'efficacité par des tests DAST et une analyse de l'observabilité.
\end{enumerate}

\section*{Structure du rapport}

Ce rapport s'articule en quatre chapitres complémentaires, encadrés par une
introduction et une conclusion générales :

\begin{itemize}
  \item Le \textbf{Chapitre 1} pose les fondements théoriques du DevSecOps et
    présente les principaux types de tests de sécurité (SAST, DAST, IAST, RASP),
    avant d'approfondir les enjeux spécifiques du DAST dans les architectures
    microservices.

  \item Le \textbf{Chapitre 2} est consacré au Service Mesh : son architecture
    bicouche (data plane / control plane), ses composants clés, et un panorama
    comparatif des solutions leaders du marché. Il examine également les apports
    concrets du Service Mesh dans une démarche DevSecOps.

  \item Le \textbf{Chapitre 3} présente l'analyse de menaces selon la méthodologie
    STRIDE appliquée à notre architecture cible, en comparant les scénarios
    d'attaque avant et après déploiement du Service Mesh.

  \item Le \textbf{Chapitre 4} décrit l'implémentation du prototype, les
    configurations de sécurité mises en place, les résultats des tests de
    validation, ainsi que le tableau de bord d'observabilité.
\end{itemize}
